\documentclass[10pt,a4paper]{amsart}
\usepackage[margin=19mm]{geometry}
\usepackage{amsmath,amssymb,mathtools}
\usepackage[scr=rsfs,cal=euler]{mathalfa}
\usepackage{xfrac}
\usepackage[unicode,hidelinks,bookmarks=false]{hyperref}
\newcommand{\ie}{i.e.\ }
\newcommand{\cf}{cf.\ }
\newcommand{\resp}{respectively\ }
\newcommand{\Subsection}{\subsection}
\providecommand{\nobreakdash}{\nobreak\hbox{-}}
\setlength{\emergencystretch}{2em}
\multlinegap0pt
\newcommand{\Q}{\mathbb{Q}}
\usepackage{amssymb}
\usepackage{enumerate}
\usepackage{xcolor}
\newtheorem{lemma}{Lemma}
\title{The Mathieu group $M_{23}$ is a Galois group over $\Q$}
\author{Xiaoyu Huang, Blake Jackson, Kyu-Hwan Lee, Bjorn Poonen, Rachel Pries, Shaowu Zhang}
\date{Source: arXiv:2608.08538v1 (CC BY 4.0)}
\begin{document}
\maketitle

\noindent\emph{Source and licence.} The Mathieu group $M_{23}$ is a Galois group over $\Q$. Version arXiv:2608.08538v1 (CC BY 4.0), distributed by arXiv under the Creative Commons Attribution 4.0 International licence, http://creativecommons.org/licenses/by/4.0/, as recorded in the arXiv licence metadata for this exact version and shown on the abstract page https://arxiv.org/abs/2608.08538v1. Full licence notice: \texttt{licenses/arxiv-2608.08538-CC-BY-4.0.txt}.\par\smallskip

\noindent\textit{Editorial scope.} Counted unit: the article's lemma lemma (source0.tex lines 456--458). The article's complete original proof of it is delivered with it (source0.tex lines 460--465, one \texttt{proof} environment of 6 lines). No mathematics is written, repaired or re-derived: the statement and the proof are the article's own lines, in the article's own notation, and no retained line is substituted or re-flowed.  No further source-local context is needed: every cross-reference of every retained line resolves inside the two delivered spans.  Nothing was omitted from any retained span, so the package carries no omission record.  No retained line contains a citation, so the wrapper carries no bibliography and no bibliography entry.  No retained line contains a figure environment, an image-inclusion command or drawing code, so no image is delivered and the package declares no asset dependency.  Editorial formatting only, in addition: an AMS \texttt{amsart} preamble that restates the article's own declarations and macros used by the retained lines, namely the environments \texttt{lemma} on separate counters, together with the standard packages those lines need.  The retained environments print in this document as Lemma 1 in order of appearance, whereas the article prints the same items with the numbering of its own full document.  The complete native source of the article as distributed in the arXiv e-print for this exact version is bundled unchanged as \texttt{sources/207172/source0.tex}.
\begin{lemma}
  The points $a,b,c$ are the vertices of a hyperbolic triangle in $D$ with corresponding interior angles $\pi/2$, $\theta$, $\theta$. 
\end{lemma}

\begin{proof}
Note that $|a|^2 < 1$ and $|c|^2 < 1$, so $a,b,c \in D$.
Since $b=0$, we have $\angle b = \arg(c) - \arg(a) = \theta$.
The hyperbolic law of cosines gives $\cosh d(a,c) = \cot \theta$.
A formula for the cosines of the angles in terms of the side lengths of a hyperbolic triangle gives $\cos(\angle a) = 0$ and $\cos(\angle c) = \cos \theta$, so $\angle a = \pi/2$ and $\angle c = \theta$.
\end{proof}

\end{document}
